\chapter{计算机辅助数控编程（APT与CAM）}
\label{ch:nc}
\sourcepage{APT补充讲义}{8、13--20}

\section{从数控程序到机床运动}
数控（NC）利用数字化指令控制机床运动和加工过程；计算机数控（CNC）由专用计算机实现相关控制功能；DNC强调计算机与多台数控设备之间的程序传送、连接或集中管理。复习时先区分\textbf{程序编制}与\textbf{程序执行}：前者决定路径和工艺，后者将指令变成机床各轴的运动。

\dfn{数控编程}{依据零件图样或模型、技术要求和工艺方案，确定加工轨迹及参数，生成符合指定机床和控制系统要求的加工程序，并进行校核。自动编程是部分或全部工作由计算机辅助完成，不表示不再需要工艺判断。}

数控加工适于复杂形状及多品种、小批量等场景，能提高重复加工的一致性、减轻操作负担并便于数字化管理。但加工精度仍受机床几何误差、热变形、切削力、装夹和控制误差等影响。“程序正确”不等于“零件一定合格”。

\subsection{一条完整的数据与控制链}
\begin{enumerate}
\item \textbf{CAD与工艺规划：}建立零件模型，确定加工特征、装夹、刀具、余量和加工顺序。
\item \textbf{CAM路径计算：}按走刀方式与误差要求计算刀位轨迹，可由线段、圆弧或样条等组成。
\item \textbf{后置处理：}结合机床结构和控制器格式，将刀位及工艺信息变为目标NC程序。
\item \textbf{CNC解释、规划与插补：}解释程序，进行路径光顺和速度规划，计算每个插补周期内各轴的指令增量或目标位置。
\item \textbf{伺服执行与反馈：}驱动电机，利用位置等反馈修正跟踪误差，形成刀具相对工件的实际运动。
\end{enumerate}
\sourcepage{APT补充讲义}{14、30--32}

\nt{后置处理发生在编程链中，插补发生在数控执行链中。后置处理可以选择或输出控制器支持的直线、圆弧等指令；控制器在运行时再把这些指令离散为周期运动命令。两者不能视为同一过程。}

\section{编程内容、坐标计算与校核}
\sourcepage{APT补充讲义}{16--20}

\subsection{六个步骤}
\begin{enumerate}
\item \textbf{工艺方案分析：}分析形状、技术要求、批量与设备，确定毛坯、工序划分和机床。
\item \textbf{工序详细设计：}确定定位夹紧、工步、刀具、切削用量、粗精加工与进退刀路线。
\item \textbf{轨迹坐标计算：}建立工件坐标系，计算基点、逼近节点及辅助运动坐标。
\item \textbf{编写加工程序：}按目标系统代码、格式与功能要求编写，并准备必要的调整和工艺文件。
\item \textbf{准备与传输程序：}将程序可靠地送入目标控制系统。讲义中的穿孔带属于历史载体示例，复习重点是程序信息的正确传递。
\item \textbf{校核与试切：}检查轨迹、行程、干涉及工艺条件，通过适当的仿真、验证和首件检查确认结果。
\end{enumerate}

\subsection{基点、节点、辅助计算}
\begin{center}
\begin{tabularx}{\linewidth}{@{}p{.18\linewidth}XX@{}}
\toprule
对象 & 含义 & 为什么计算\\
\midrule
基点 & 轮廓相邻几何元素的交点、切点等 & 确定线段和圆弧等原始元素的连接位置\\
节点 & 用分段直线或圆弧逼近非圆曲线时的连接点 & 使近似轨迹满足误差限\\
辅助坐标 & 引入、退出、换坐标及必要圆整等 & 保证路径衔接与程序坐标一致\\
\bottomrule
\end{tabularx}
\end{center}
刀具有尺寸，编程对象可能是零件轮廓，也可能是刀心轨迹。二者需通过刀具补偿或偏置联系。选择哪一种，取决于CAM输出方式、程序约定与控制器功能；不能将零件边界坐标直接当作刀心坐标。

\ex{圆弧离散的弓高误差（复习推导）}{取$R>0$，用弦长$0\leq\ell\leq2R$逼近相应的小圆弧。若允许弓高为$0\leq\varepsilon\leq R$，圆心到弦的距离为$\sqrt{R^2-(\ell/2)^2}$，因此最大弓高为
\[
\varepsilon=R-\sqrt{R^2-(\ell/2)^2},\qquad
\ell\leq2\sqrt{2R\varepsilon-\varepsilon^2}.
\]
当$\varepsilon\ll R$，可用$\ell\lesssim\sqrt{8R\varepsilon}$估算。若$R=50\,\mathrm{mm}$，允许弓高$0.01\,\mathrm{mm}$，则$\ell_{\max}\approx2.00\,\mathrm{mm}$。这是理想几何逼近误差，不包括刀具、装夹及机床运动误差。}

节点越密，几何近似通常越精细，但程序量与处理负担会增加，短线段还会影响速度规划。应根据曲率与容差选取步长，而不是无条件增加节点。

\section{手工、语言与图形交互式自动编程}
\sourcepage{APT补充讲义}{15、20、25--26}

\begin{center}
\begin{tabularx}{\linewidth}{@{}p{.18\linewidth}XX@{}}
\toprule
方法 & 主要过程 & 适用性与限制\\
\midrule
手工编程 & 人工分析、计算并写出NC程序 & 简单几何和较少程序段较直观；复杂空间轨迹计算负担大\\
语言自动编程 & 用APT等描述几何及运动，由系统计算刀位并后置 & 可表达复杂任务，但源程序需学习、阅读，几何不够直观\\
图形交互式 & 在CAD模型上选择几何并设置工艺，由CAM生成路径 & 便于可视化和设计制造信息连接；仍依赖正确工艺和后置配置\\
\bottomrule
\end{tabularx}
\end{center}

通用系统可按功能侧重分为\textbf{几何处理}与\textbf{工艺处理}：前者侧重形状和轨迹计算，后者还辅助加工方法和工艺参数的选择。讲义以APT和EXAPT说明这种差异。不能因为系统会自动算轨迹，就认为它已自动解决全部工艺规划。

语言自动编程系统的工作可概括为\textbf{输入翻译、轨迹计算、后置处理}三个阶段。输入翻译检查并解释源程序；轨迹计算根据几何、刀具及运动条件形成刀位数据；后置处理生成指定设备的程序。

\section{APT的几何、运动与工艺描述}
\sourcepage{APT补充讲义}{21--24}

APT（Automatically Programmed Tools）是一种自动编程语言及相应系统。理解它的重点是将“定义什么几何”“怎样运动”“采用什么工艺”分开阅读。

\subsection{符号、词汇与语句}
逗号分隔参数，斜杠分隔主要词与参数部分，也可在计算式中表示除法；星号表示乘法，双星号表示幂。常见几何词有\texttt{POINT}、\texttt{LINE}、\texttt{CIRCLE}、\texttt{PLANE}；几何关系词如\texttt{PARLEL}、\texttt{PERPTO}、\texttt{TANTO}。

\begin{quote}\ttfamily
P1 = POINT/0, 0, 0\\
P2 = POINT/20, 0, 0\\
C1 = CIRCLE/0, 0, 25\\
FROM/0, 0, 30\\
GOTO/20, 0, 30
\end{quote}
上面仅示意几何定义、起点声明与点位移动的语义。\texttt{P1}、\texttt{P2}是标识符，圆语句给出平面圆的中心与半径。完整加工程序还需明确单位、刀具、工艺、控制系统与进退刀条件；这些片段用于语言理解。

\subsection{三个控制面：PS、DS、CS}
\dfn{零件面、导向面、检查面}{零件面PS（Part Surface）用于限制刀具与待加工表面的关系，通常控制切削深度；导向面DS（Drive Surface）引导刀具运动，控制侧向关系；检查面CS（Check Surface）限定当前运动段的终止条件。刀具保持对PS、DS规定的关系，运动至满足CS条件，再切换下一运动段。}

三个面是\textbf{约束角色}，不一定都必须是平面。它们可以是由几何实体定义的控制表面。只给一条终止面而不说明刀具与其余面的关系，通常不能唯一确定空间刀具运动。

\subsection{连续运动语句的读法}
初始连续运动的示意格式为
\begin{quote}\ttfamily
GO/TO, DS, TO, PS, PAST, CS
\end{quote}
三个控制面的书写顺序是\textbf{DS、PS、CS}，各面前的限定词描述刀具与该面的目标关系。上例先说明与导向面、零件面的关系，再说明经过检查面后停止。它是语法模板，实际程序应以已定义的几何实体代替DS、PS、CS。

后续轮廓运动可写成例如
\begin{quote}\ttfamily
GOFWD/DS, TO, CS
\end{quote}
表示沿当前前进方向、以DS导向，至满足对CS的TO条件结束；也可按几何关系指定\texttt{TANTO}等相切终止条件。后续段继承的零件面关系及各几何定义也必须明确，不能只读方向词就确定刀轨。

\subsection{读懂TO、ON、PAST}
沿当前运动方向观察，到检查面的停止条件不同：\texttt{TO}在到达前侧相切位置停止；\texttt{ON}使刀具中心到面上；\texttt{PAST}使刀具经过面后在另一侧相切位置停止。\texttt{FROM}指定初始刀位，\texttt{GOTO}指定目标刀位；连续轮廓运动还可使用\texttt{GOFWD}、\texttt{GORGT}、\texttt{GOLFT}等方向词。

\ex{三个停止条件为什么会给出不同刀位？}{在二维截面内，设圆形刀具半径$r$，中心沿$+x$移动，检查直线为$x=c$。在可达到且无其他干涉的条件下，\texttt{TO}、\texttt{ON}、\texttt{PAST}对应中心位置分别为$c-r$、$c$、$c+r$。若$c=50\,\mathrm{mm}$、$r=5\,\mathrm{mm}$，则为45、50、55 mm。三个数指刀具中心，而非同一个刀刃接触点。}

\nt{原课件第22页个别英文单词存在拼写混排，如PASS与PAST、GOLET与GOLFT。这里采用通常的PAST、GOLFT写法；理解考试题时应根据语义识别。方向词依赖当前刀具运动方向，不能简单等同于固定世界坐标轴的正负方向。}

\reviewfig[.46]{ch09-apt-profile.pdf}{由直线与圆弧构成的APT轮廓实例（APT讲义，PDF第24页局部）。}{fig:nc-apt-profile}

图\ref{fig:nc-apt-profile}的读法是先识别圆弧$C_1$至$C_4$和直线$L_1$至$L_3$，再检查相切关系与走刀衔接。定义一组轮廓元素只建立几何约束；还要说明刀具在哪一侧、按什么方向运动、在哪个检查面结束，才能形成确定刀轨。

\section{后置处理：刀位数据如何适配机床}
\sourcepage{APT补充讲义}{23、26、30}

\dfn{CLDATA与后置处理}{CLDATA（Cutter Location Data）是表示刀具位置、方向及相关加工信息的中间刀位数据。后置处理器结合目标机床的结构、坐标约定、轴运动能力与控制系统功能，将中间数据转化为设备可接受的NC程序。}

后置处理需完成或协调机床坐标转换、运动轴解算、程序格式、准备功能与辅助功能、速度及相关限制检查，并输出必要的加工和操作信息。对多轴机床，同一刀位与刀轴方向还可能对应不同轴角解，需要兼顾行程、连续性和机床运动学。

\ex{更换机床时，为什么不只改程序后缀？}{同一零件刀轨换到另一机床，可能遇到坐标原点、轴配置、旋转轴正方向、代码格式或辅助功能差异。刀位计算结果可在条件合适时复用，但要使用相匹配的后置处理，并重新检查行程、干涉、刀具与工艺条件。刀位数据可移植，不等于成品NC程序可任意移植。}

\subsection{CAD/CAM集成的几种路径}
可以由人员依据图样把设计转成NC编程输入；可以让CAD与CAM共享软件内部模型；也可通过接口交换几何和产品数据。若进一步接入CAPP，几何之外还能传递特征、加工顺序和参数等信息。应同时考虑几何一致、工艺语义一致和变更后的版本同步。

\section{精度、效率与验证}
\sourcepage{APT补充讲义}{18、30--35}

插补将路径变成周期指令：既可理解为对离散点之间的路径加密，也可理解为把连续曲线离散为控制周期位置。速度规划还应满足几何容差和速度、加速度、跃度及各轴能力等约束。路径拐角很尖时，若不允许改变几何，就必须相应减速；若要保持较高速度，则可能需要在容差内光顺。

\begin{center}
\begin{tabularx}{\linewidth}{@{}p{.18\linewidth}XX@{}}
\toprule
误差类别 & 典型来源 & 对应思路\\
\midrule
几何误差 & 零件形状偏差、装配失准、轴运动误差 & 结构与装配控制、标定及几何补偿\\
热误差 & 主轴和丝杠热伸长、切削热传递 & 热设计、热稳定、温度与位移补偿\\
力误差 & 切削力、自重与装夹变形 & 提高适当刚度、合理支撑与工艺参数\\
控制误差 & 插补近似、跟踪滞后、轴间不匹配 & 合理插补规划、伺服调节与轮廓控制\\
\bottomrule
\end{tabularx}
\end{center}

提高刚度有助于减小受力变形，但运动部件过重会限制高加速度，增加驱动负担；精度与效率需要结合结构和工艺共同考虑。薄壁件、长悬伸件还应注意变形与振动，不能只依据理想几何模型决定进给速度。

\nt{轨迹仿真、空运行与首件检测回答不同问题：仿真检查模型中可表示的运动和干涉，空运行可核对实际机床的运动关系，首件检测检验综合制造结果。任何一种检查都不能自动覆盖其余检查的全部目的。}

\Needspace{11cm}
\section{复习自测}
\begin{enumerate}
\item 数控编程与插补的输出有何不同？\textit{答：前者输出加工程序，后者输出各控制周期的轴运动指令。}
\item 非圆曲线的节点数由什么决定？\textit{答：曲线形状与曲率、逼近方法和允许误差等。}
\item 零件轮廓能否直接作为铣刀刀心轨迹？\textit{答：通常需考虑刀具半径及偏置或补偿约定。}
\item PS、DS、CS分别约束什么？\textit{答：深度或主接触关系、导向关系、当前段终止条件。}
\item \texttt{FROM}是否就是从当前位置移动到指定点？\textit{答：它声明初始刀位；目标移动用相应运动语句描述。}
\item 为什么后置处理与机床相关？\textit{答：机床运动学、坐标、功能和程序格式不同。}
\item CLDATA与G代码是否同义？\textit{答：前者是中间刀位数据，后者是目标控制程序中的代码表达。}
\item 空运行正确后首件仍超差，有哪些原因？\textit{答：刀具补偿、受力变形、热误差、装夹或测量等可能尚未被验证。}
\item 为什么增加路径节点不一定提高实际加工效率？\textit{答：程序增大、短段与拐角增多，会增加计算和速度规划负担。}
\end{enumerate}
